% Title page, summary, contents and lists of figures and tables.

\begin{titlepage}
  \thispagestyle{empty}
  \vspace*{30mm}
  {\sffamily\bfseries\fontsize{40}{44}\selectfont\color{accent}Yap\par}
  \vspace{6mm}
  {\sffamily\bfseries\LARGE Technical report\par}
  \vspace{4mm}
  {\sffamily\large\color{muted}Architecture, data flow and features of the Yap mobile app\\and its Supabase backend\par}
  \vspace{10mm}
  {\color{accent}\rule{28mm}{1.2pt}\par}
  \vspace{10mm}
  {\small
  \begin{tabular}{@{}L{34mm}L{100mm}@{}}
    \textsf{\textbf{Subject}} & The YapMap repository, branch \code{main} after pull request 74 and the migration renames (October 2026) \\
    \textsf{\textbf{Platforms}} & iOS and Android (Expo SDK 57, React Native 0.86) \\
    \textsf{\textbf{Backend}} & Supabase: Postgres with PostGIS, Auth, Realtime, Storage and \code{pg_cron} \\
    \textsf{\textbf{Course}} & COMP 2160 (Section 01), Mobile App Development 1, final project \\
    \textsf{\textbf{Team}} & Samar Khajuria and Manik Singh \\
    \textsf{\textbf{Date}} & October 2026 \\
  \end{tabular}\par}
  \vfill
  {\footnotesize\color{muted}
  Copyright \textcopyright{} 2026 Samar Khajuria. All rights reserved. The software described here and its
  source code are proprietary and may not be used, copied, modified or distributed without the written
  permission of the copyright holder.\par}
\end{titlepage}
\pagenumbering{roman}

\chapter*{Summary}
\addcontentsline{toc}{chapter}{Summary}
\markboth{Summary}{}

Yap is a mobile app that helps people who share a physical space, such as a café or a campus lounge,
start a conversation in person. A user goes live for 30, 60 or 120 minutes, sees the other live users
within about 100 metres, and sends one of them a request to talk. If the request is accepted within five
minutes, the two users enter a private chat that exists to help them find each other and meet. Users can
also place small public events on a map, which anyone within 5 kilometres can join, each with its own
group chat.

The client is a single Expo and React Native code base written in TypeScript. Expo Router maps files to
screens; one store in \code{src/state/store.ts} holds all shared state and owns every Realtime
subscription; and one module, \code{src/data/api.ts}, makes every call to the backend. The backend is a
Supabase project with no application server. Postgres row-level security denies direct access by default,
and every write other than the user's own profile columns, an unblock and the user's own photo files
goes through a \code{SECURITY DEFINER} database function that derives the caller from the session and
checks what that action needs: identity, state, proximity, blocks or rate limits. PostGIS answers the proximity
questions, Realtime delivers requests and messages, a private Storage bucket holds photos behind
short-lived signed URLs, and a \code{pg_cron} job ends expired sessions, requests and stale chats every
minute.

Privacy shapes the design throughout. Live positions are written once per session, are never returned to
any client, and are deleted the moment a user goes offline. Other users see a person card that carries a
first name and last initial but no location or distance. A photo can be loaded only by someone who can
currently see its owner. The project accepted one deliberate trade-off to keep the promise that a location
is deleted on going offline: the server cannot rate-limit the distance between two separate live
sessions, so it limits how often a location may change instead.

Safety follows the app stores' rules for apps where people post and meet. Every user agrees to community
guidelines at sign-up, can report and block from every screen that shows another person, and cannot post
slurs or abusive words, which the database refuses. The owner reviews reports at least daily and can
suspend an account. Continuous integration checks the app, the migration files and a database rebuilt
from every migration on each pull request, and store builds are made with EAS Build.

This report describes the architecture (\cref{ch:architecture}), the data model (\cref{ch:data-model})
and the server interface (\cref{ch:server-interface}); traces each major flow with sequence diagrams
(\cref{ch:flows}); catalogues every feature with the screens, functions and server rules that implement
it (\cref{ch:features}); and covers security and privacy, reliability, performance, configuration,
validation, and known limitations in the chapters that follow.

\clearpage
{\hypersetup{linkcolor=ink}\tableofcontents}
\clearpage
\listoffigures
\listoftables
